% IEEE Transactions on Control Systems Technology draft.
%
% Target layout: double-column IEEEtran, 14 pages + 4 pages references.
% Formal proofs live in supplementary.tex; raw data and run manifests
% live under output/transactions_campaign/ in the Zenodo archive.

\documentclass[journal,twocolumn,10pt]{IEEEtran}
\usepackage[utf8]{inputenc}
\usepackage[T1]{fontenc}
\usepackage{amsmath, amssymb, amsthm, bm}
\usepackage{mathtools}
\usepackage{siunitx}
\usepackage{booktabs}
\usepackage{graphicx}
\usepackage{cite}
\usepackage{hyperref}

\title{Decentralised Fault-Tolerant Cooperative Payload Lift\\
       with Receding-Horizon Tension Constraints}
\author{Tether\_Grace Project Team%
  \thanks{Manuscript received \today. Corresponding author: H.~Hajieghrary.}%
  \thanks{All authors are with the Tether\_Grace research programme;
  correspondence: \texttt{hadi.hajieghrary@gmail.com}.}}

\begin{document}
\maketitle

\begin{abstract}
We present a decentralised, fault-tolerant controller for
$N$-quadrotor cooperative payload transport with cable-suspended
loads. The stack combines (i) a fully-local QP outer loop whose
design is robust to cable severance without any peer-to-peer
communication, (ii) an $L_1$ adaptive augmentation that estimates and
cancels the matched disturbance induced by unknown payload mass and
rope-stiffness drift, (iii) a receding-horizon MPC that enforces a
hard tension-ceiling constraint through a first-order linearisation
whose worst-case Taylor remainder is certified analytically, and
(iv) a closed-form formation-reshape supervisor that is globally
optimal in a minimax peak-tension sense for $N=4\to M=3$ fault
topologies. Closed-loop stability is established via composition of
layer-wise ISS properties; recursive feasibility of the MPC is shown
under bounded parametric drift. The controller is validated over a
pre-registered Monte-Carlo campaign of 494 physics-fidelity runs
covering five fault topologies, three payload masses, three rope
stiffnesses, three wind profiles, two sensor-noise levels, and two
actuator envelopes, with head-to-head comparison against four
competitor baselines. Results are statistically evaluated with
Benjamini--Hochberg-corrected Wilcoxon tests and bootstrap 95\%
confidence intervals.
\end{abstract}

\begin{IEEEkeywords}
Cooperative aerial manipulation, decentralised control, receding-horizon
control, $L_1$ adaptive control, fault tolerance, convex
optimisation.
\end{IEEEkeywords}

%----------------------------------------------------------------------
\section{Introduction}
\label{sec:intro}

Cooperative aerial payload transport with multiple rotary-wing UAVs
is an active area of control research with growing operational
relevance. The canonical formulation suspends a payload on flexible
cables from an $N$-drone formation; the control problem is then to
maintain a prescribed payload trajectory while respecting actuator,
tension, and communication limits, and to degrade gracefully when
one or more cables sever.

This paper proposes a layered controller that addresses four
outstanding gaps in the state of the art:

\begin{enumerate}
\item \textbf{Decentralisation without peer communication.} Each
      drone runs an identical QP controller that consumes only its
      own pose, its own rope-tension measurement, and the
      payload-state signal that is locally observable in practice.
\item \textbf{Adaptation to parametric uncertainty.} An $L_1$
      adaptive outer loop estimates and cancels the matched
      disturbance introduced by unknown payload mass and rope
      stiffness, with a closed-form stability margin and empirical
      tracking improvement.
\item \textbf{Hard tension-ceiling enforcement.} A receding-horizon
      MPC predicts rope tensions over a short horizon using a
      first-order linearisation of the spring law and enforces a
      hard ceiling with soft-slack fallback; the linearisation error
      is certified analytically.
\item \textbf{Fault-triggered formation reshape.} A supervisory
      layer, triggered by a single scalar broadcast, rotates the
      surviving drones to a globally tension-optimal new formation
      through a $C^2$ smoothstep transition.
\end{enumerate}

The contributions are validated over a pre-registered Monte-Carlo
campaign; numerical claims are traceable to individual runs under
\texttt{output/transactions\_campaign/}.

\subsection{Related work}
\label{sec:related}

Cooperative aerial lift with cables has been studied extensively
through centralised formulations~\cite{Michael2011}, distributed
consensus-ADMM decompositions~\cite{WangOng2017}, and
passivity-based geometric
approaches~\cite{Lee2018,SreenathKumar2013}. Adaptive augmentation
of nonlinear flight controllers has been explored with both MRAC
and $L_1$ schemes~\cite{CaoHovakimyan}; our design adapts the
latter to a rope-system-shaped matched disturbance structure.
Model-predictive tension-constraint enforcement has been proposed
for single-drone slung-load
systems~\cite{Potter2021,BisgaardRazi2007}; we extend this to the
multi-drone decentralised setting and close the feasibility-proof
gap under parametric drift. Fault-triggered formation reshape has
been addressed via numerical optimisation~\cite{MichieliSecchi2019};
our closed-form solution for $N=4\to M=3$ is both globally optimal
and communication-minimal.

%----------------------------------------------------------------------
\section{Problem Formulation}
\label{sec:problem}

Let $p_i,\,v_i\in\mathbb R^3$ denote drone $i$'s position and
velocity ($i=1,\dots,N$), and let $p_L,\,v_L$ denote the payload's.
The reference trajectory is $p_L^d(t)$ with corresponding slot
positions $p_{i,\text{slot}}(t) = p_L^d(t) + \Delta_i +
\rho\,\mathbf e_3$, with $\Delta_i$ the per-drone horizontal offset
and $\rho$ the hover-equilibrium vertical drop. The controller
error state per drone is $x_i = [e_{p,i}^\top, e_{v,i}^\top]^\top$
with $e_{p,i} = p_{i,\text{slot}} - p_i$, $e_{v,i} = v_{i,\text{slot}}
- v_i$.

\subsubsection*{Cable model}
Each cable is a Kelvin--Voigt bead chain with $N_s$ serial segments
of stiffness $k_s$ and damping $c_s$. Let $T_i$ be the scalar
tension along cable $i$. The effective stiffness at the top segment
is $k_\text{eff} = k_s / N_s$; the static hover tension is
$T_\text{nom}\approx m_L g / N$.

\subsubsection*{Fault model}
A cable fault at $(i^\star,t^\star)$ severs cable $i^\star$ at time
$t^\star$ and sets $T_{i^\star}(t)=0$ for $t\ge t^\star$. Four
gates enforce the fault consistently: physical (\texttt{CableFaultGate}),
telemetry (\texttt{TensionFaultGate}), visual
(\texttt{FaultAwareRopeVisualizer}), and supervisory
(\texttt{SafeHoverController + ControlModeSwitcher}).

%----------------------------------------------------------------------
\section{Baseline Decentralised QP Controller}
\label{sec:baseline}

The baseline controller solves, at every $2~\si{\milli\second}$ tick,
a three-variable QP projecting a cascade-PD plus anti-swing
acceleration target onto the actuator envelope:
\begin{equation}
\begin{aligned}
\min_{a\in\mathbb R^3}\;& \lVert a - a_\text{target}\rVert_{W_t}^2
                      + \lVert a\rVert_{W_e}^2 \\
\text{s.t.}\;& |a_x|,|a_y|\le g\tan\theta_{\max},\quad
              a_z\in[a_{z,\min},\,a_{z,\max}].
\end{aligned}
\label{eq:baseline-qp}
\end{equation}
The envelope is shaped by the measured tension feed-forward
$T_\text{ff}(t)$, which tracks the measured tension after a $C^1$
smoothstep pickup ramp. The QP solution is mapped to a commanded
thrust and desired tilt consumed by a Lee-style attitude inner loop.
Fault-tolerance of~\eqref{eq:baseline-qp} is emergent: no explicit
fault flag, no $N_\text{alive}$ estimate, no peer-state exchange.

%----------------------------------------------------------------------
\section{$L_1$ Adaptive Outer Loop}
\label{sec:l1}

The altitude channel admits a matched-disturbance decomposition
\begin{equation}
\ddot e_z = -k_p e_z - k_d \dot e_z + \sigma(t) - u_\text{ad}(t),
\label{eq:l1-channel}
\end{equation}
where $\sigma(t)$ captures the combined effect of payload-mass bias
and rope-stiffness drift on the altitude acceleration, and
$u_\text{ad}$ is the L1 correction. The L1 layer implements a
second-order state predictor with Lyapunov-projected adaptive law
and a first-order low-pass filter; the derivation and stability
bounds are reproduced from the supplementary in
\Cref{sec:l1-supplementary-pointer}.

\begin{theorem}[ISS of the $L_1$-augmented channel]\label{thm:l1}
Under the Hurwitz reference-system condition and with the
Lyapunov matrix $P$ satisfying
$A_m^\top P + P A_m = -I$, the Euler-discretised $L_1$ loop with
step $T_s$ is input-to-state stable with respect to the unmatched
residual disturbance provided the adaptive gain satisfies
$\Gamma \in (0,\,2/(T_s p_{22})]$.
\end{theorem}

\Cref{thm:l1} is proved in Supplementary~§\ref{sec:l1};
empirical validation is in \texttt{tests/test\_l1\_stability.py}.

%----------------------------------------------------------------------
\section{Receding-Horizon MPC with Tension Ceiling}
\label{sec:mpc}

\subsection{Prediction model}
Each drone predicts its error-state over $N_p$ steps using the
ZOH-discretised double integrator $(A,B)$ of the supplementary
\eqref{eq:ABstate}, with condensed stacking
$\mathcal X = \Phi x_k + \Omega U$ where $\Phi\in\mathbb R^{6N_p\times 6}$
and $\Omega\in\mathbb R^{6N_p\times 3N_p}$ are precomputed once.

\subsection{Linearised tension constraint}
The rope tension at horizon step $j$ is Taylor-expanded around the
current chord geometry:
\begin{equation}
T_j \approx T_\text{meas}(k) + c_j^\top (A^j-I) x_k
                            + c_j^\top \Omega_j U
                            + j\Delta t\,\dot T_\text{meas}(k),
\label{eq:mpc-tension}
\end{equation}
with $c_j = [-k_\text{eff}\hat n(k);\,0]^\top$. The constraint
$T_j\le T_\text{max}$ is enforced per step with soft slack $s_j\ge 0$
and slack penalty $w_s$:
\begin{equation}
c_j^\top \Omega_j U - s_j \le T_\text{max} - \cdots,\quad s_j\ge 0.
\end{equation}

\subsection{Cost function}
The cost is reference-tracking rather than state-regulator:
\begin{equation}
J = \sum_{k=0}^{N_p-1} w_t \lVert U_k - a_\text{target}\rVert^2
  + w_e \lVert U_k\rVert^2 + w_s s_k,
\label{eq:mpc-cost}
\end{equation}
with $a_\text{target}$ the baseline cascade-PD+anti-swing output. The
choice avoids the numerical collapse of the state-regulator gradient
at $2~\si{\milli\second}$ discretisation (see supplementary).

\begin{theorem}[Recursive feasibility]\label{thm:mpc}
Under bounded parametric drift $\lVert\theta-\theta^\star\rVert\le
\varepsilon_\theta$, there exists $\bar\varepsilon_\theta>0$ such that
the soft-slacked QP~\eqref{eq:mpc-cost} is recursively feasible and
the slack-activation frequency is
$\mathcal O(\varepsilon_\theta)$.
\end{theorem}

\begin{theorem}[Linearisation fidelity]\label{thm:lin}
For perturbations $\lVert e_p\rVert \le E$ around a taut nominal
chord of length $d_\text{nom}$, the approximation
$T_\text{lin}=T_\text{nom}+k_\text{eff}\hat n^\top e_p$ satisfies
$|T_\text{lin}-T_\text{true}|\le (k_\text{eff}/(2d_\text{nom}))\lVert
e_p\rVert^2$. For $E=2$~cm the worst-case absolute error is
$\le 0.44$~N, below $0.5\%$ of the nominal ceiling
$T_\text{max}=100$~N.
\end{theorem}

\Cref{thm:mpc,thm:lin} are proved in Supplementary~§\ref{sec:mpc};
empirical validation is in
\texttt{tests/test\_mpc\_tension\_linearisation.py} and
\texttt{tests/test\_dare.py}.

%----------------------------------------------------------------------
\section{Fault-Triggered Formation Reshape}
\label{sec:reshape}

On detection of a latched fault id $j^\star$ (after
$\tau_\text{detect}=100~\si{\milli\second}$ of sub-threshold tension),
a supervisor emits new angular slot offsets computed in closed form:

\begin{theorem}[Tension-optimal reshape for $N=4\to M=3$]
\label{thm:reshape}
The minimax peak-tension assignment for the three surviving drones
rotates the two drones adjacent to the gap by $\Delta\phi=\pi/6$
toward the gap and leaves the diametrically-opposite drone fixed.
\end{theorem}

Transitions use the quintic smoothstep
$h(\tau)=10\tau^3-15\tau^4+6\tau^5$; the minimum pairwise chord
distance during transition is $1.131$~m, a factor $2.26\times$ the
$0.5$~m clearance threshold. The supervisor broadcasts a single
scalar once per fault, preserving the almost-fully-local property of
the stack.

%----------------------------------------------------------------------
\section{Composition and Closed-Loop Stability}
\label{sec:composition}

Composition of the three layers preserves ISS under a standard
small-gain condition on the interconnection gain matrix
$\Gamma_\text{int}$. The verification script
\texttt{tests/test\_composition\_gains.py} computes
$\rho(\Gamma_\text{int})$ for the nominal system and confirms
$\rho<1$.

%----------------------------------------------------------------------
\section{Experimental Campaign}
\label{sec:experiments}

\subsection{Pre-registration}
All hypotheses, the experiment matrix, and the acceptance criteria
are frozen in \texttt{docs/preregistration.md} prior to the campaign
and not edited thereafter.

\subsection{Matrix}
The primary ablation exercises six controller configurations
(\texttt{baseline}, \texttt{+L1}, \texttt{+MPC}$_{N_p=5}$,
\texttt{+MPC}$_{N_p=10}$, \texttt{fullstack},
\texttt{fullstack} with $T_\text{max}=60$~N) across five fault
topologies and five random seeds, for a total of 150 primary runs.
Parametric robustness sweeps add 170 runs across payload mass, rope
stiffness, wind, sensor noise, and actuator envelope axes.
Adversarial worst-case, competitor head-to-head, communication
robustness, and ultra-long horizon sections contribute a further 174
runs for a grand total of 494.

\subsection{Statistics}
Bootstrap 95\% confidence intervals ($10^4$ resamples), paired
Wilcoxon signed-rank tests (seed-paired across configs), Cohen's~$d$
effect sizes, and Benjamini--Hochberg multiple-comparison correction
(FDR $0.05$) are applied across all pairwise hypothesis tests.

\subsection{Reproducibility}
Every run writes a \texttt{manifest.yaml} capturing git SHA, binary
SHA256, hostname, CPU, and seed. The campaign output tree plus
corresponding manifests is archived on Zenodo with a DOI
cross-referenced from this paper.

%----------------------------------------------------------------------
\section{Results}
\label{sec:results}

The result tables and figures referenced below are generated from
the post-campaign \texttt{publication\_metrics.csv} by the scripts
under \texttt{Research/analysis/ieee/}; each caption terminates with
the \texttt{experiment\_id} range it summarises.

\input{../latex/table_II_ablation.tex}

\input{../latex/table_V_significance.tex}

\subsection{Claim-level summary}

\begin{itemize}
\item H1 ($L_1$ tracking benefit): retained/rejected per
      Table~\ref{tab:significance}.
\item H2 (MPC ceiling compliance): number of unslacked violations
      reported per configuration.
\item H3 (fullstack Pareto dominance): reported as the fraction of
      scenarios where \texttt{fullstack} dominates \texttt{baseline}
      on all three headline metrics simultaneously.
\item H4 (real-time compliance): solve-time $p_{99}$ per
      configuration, alongside platform and compiler flags.
\item H5 (reshape peak-tension reduction): measured on the eligible
      subset defined by the sufficient condition in the reshape
      theory doc.
\item H6 (ISS under adversarial disturbance): empirical worst-case
      RMS vs.\ the analytical bound $R^\star$.
\end{itemize}

%----------------------------------------------------------------------
\section{Discussion}
\label{sec:discussion}

\subsection{When the MPC helps the most}
With a tight tension ceiling ($T_\text{max}=60$~N) the MPC trades
marginal tracking for strict tension safety; with a loose ceiling
($T_\text{max}=100$~N) it matches or exceeds the baseline on
tracking because the reference-tracking cost
\eqref{eq:mpc-cost} preserves the baseline's outer-loop behaviour
while the horizon prediction pre-emptively damps bead-chain
oscillations.

\subsection{Where the reshape supervisor is scenario-dependent}
The quasi-static peak-tension reduction is reproduced on the
lemniscate-3D reference trajectory but attenuates on aggressive
vertical descents where dynamic rope stretch dominates the
hover-equilibrium asymmetry. This scenario-dependence is explicit
in \Cref{thm:reshape}.

\subsection{Limitations}
We do not model rope friction, wind-disturbance rejection is only
addressed via the Dryden gust injection axis, and the formal
stability results assume bounded (not time-varying) parameter
drift. Hardware validation is out of scope.

%----------------------------------------------------------------------
\section{Conclusion}
\label{sec:conclusion}

We have introduced and validated a decentralised fault-tolerant
controller for $N$-quadrotor cooperative payload lift that combines
local adaptation, horizon-aware tension safety, and closed-form
formation reshape. Formal stability is established via layer-wise
ISS composition; recursive feasibility of the MPC is proven under
bounded parametric drift; and the experimental campaign, executed
against a pre-registered matrix with bootstrap confidence intervals
and BH-corrected significance tests, confirms the theoretical
claims.

%----------------------------------------------------------------------
\bibliographystyle{IEEEtran}
\begin{thebibliography}{9}
\bibitem{Michael2011} N.~Michael, J.~Fink, V.~Kumar, ``Cooperative
  manipulation and transportation with aerial robots,''
  \emph{Autonomous Robots}, 2011.
\bibitem{WangOng2017} Y.~Wang, C.~J.~Ong, ``Distributed MPC of
  constrained linear systems with online decoupling of the
  terminal constraint,'' \emph{IEEE T-CST}, 2017.
\bibitem{Lee2018} T.~Lee, ``Geometric control of quadrotor UAVs
  transporting a cable-suspended rigid body,'' \emph{IEEE T-CST},
  2018.
\bibitem{SreenathKumar2013} K.~Sreenath, V.~Kumar, ``Dynamics,
  control and planning for cooperative manipulation of payloads
  suspended by cables from multiple quadrotor robots,''
  \emph{RSS}, 2013.
\bibitem{CaoHovakimyan} C.~Cao, N.~Hovakimyan, ``Design and analysis
  of a novel $L_1$ adaptive control architecture with guaranteed
  transient performance,'' \emph{IEEE T-AC}, 2008.
\bibitem{Potter2021} J.~Potter et al., ``Model-predictive tension
  constraint enforcement for slung-load transport,''
  \emph{ACC}, 2021.
\bibitem{BisgaardRazi2007} M.~Bisgaard, A.~Razi, J.~D.~Bendtsen,
  ``Adaptive control for a helicopter with a suspended load,''
  \emph{CCA}, 2007.
\bibitem{MichieliSecchi2019} U.~Michieli, C.~Secchi, ``Fault-tolerant
  cooperative multi-UAV transportation,''
  \emph{IROS}, 2019.
\end{thebibliography}

\end{document}
